\documentclass[12pt]{article}
\usepackage[utf8]{inputenc}
\usepackage{amsmath}
\usepackage{mathrsfs}
\usepackage{amssymb}
\usepackage{amsfonts}

\DeclareMathOperator{\logit}{logit}
\DeclareMathOperator*{\argmax}{arg\,max}
\DeclareMathOperator*{\argmin}{arg\,min}

\usepackage[
    a4paper,
    left = 2.5cm,
    right = 2.5cm,
    top = 2.5cm,
    bottom = 2.5cm
]{geometry}

\usepackage{graphicx}
\usepackage[headings]{fullpage}
\usepackage{hyperref}
\usepackage{fancyhdr}
\usepackage{IEEEtrantools}
\usepackage{xcolor}

%\usepackage[
    %backend=biber,
    %style=apa,
    %maxcitenames=2,
    %mincitenames=1,
    %sorting=ynt
    %]{biblatex}
% \addbibresource{bibfile.bib} 
\usepackage[numbers,sort&compress]{natbib} 

\setlength{\headheight}{15pt}
\pagestyle{fancy}
\fancyhf{}
\rhead{Angel Gómez, Berk Bahar, Jingyu Zheng} 
\lhead{SBI - Three Body Problem}
\rfoot{Page \thepage}

\begin{document}


\pagestyle{empty} 

\begin{titlepage}
        \centering
        \vspace*{\baselineskip}
        
        {\huge Inferring Initial Conditions in the\\ Three Body Problem}\\[0.2\baselineskip]
        
        \vspace*{\baselineskip}
        
        {\Large --- Project Report ---\\
          Simulation Based Inference\\
          Group 33\\
          [\baselineskip]}
        \vspace*{\baselineskip}
        
        {\LARGE 
            Angel Gómez, 280844\\[10pt]
            Berk Bahar, 274028\\[10pt]
            Jingyu Zheng, 274006\\[10pt]
        }
        
        \vfill
        August 2026 \par
        \vspace*{\baselineskip}
        {\itshape TU Dortmund University\par}
\end{titlepage}

\tableofcontents 
\clearpage       


\pagestyle{fancy}       
\pagenumbering{arabic}   






\section{Introduction}

The three-body problem is one of the most famous challenges in physics: predicting
the motion of three massive bodies interacting through gravity. Although the laws of
Newtonian gravity are simple, the behaviour of three interacting bodies is not.
Unlike a two-body system, which has a closed-form solution, a three-body system is
\emph{non-integrable} and generically \emph{chaotic}, so its long-term behaviour
becomes effectively unpredictable.
 
More technically, if two systems start from almost identical initial conditions,
their trajectories do not stay close: they separate exponentially fast, with the
separation growing as $e^{\lambda t}$, where $\lambda$ is the Lyapunov exponent.
Even the tiniest initial difference eventually leads to completely different outcomes.
 

% Figure 1 moved to appendix

Simulating the system \emph{forwards} in time is straightforward on modern computers;
the \emph{inverse} problem is much harder. In this project we focus on that inverse
problem: given only a short, noisy observation of the trajectory $\mathbf{x}$ at some
later time, can we recover the initial conditions $\theta$ at $t=0$?
 
Because chaotic dynamics erase information about the system's history, many different
starting configurations can produce almost the same observed trajectory. Reporting a
single best-guess estimate would therefore be misleading. To be statistically honest,
we report a full posterior distribution
$p(\theta \mid \mathbf{x}) \propto p(\mathbf{x} \mid \theta)\, p(\theta)$, capturing
all starting states consistent with the observation. For a chaotic system, a wide,
uncertain posterior is the correct and expected outcome, and the failure mode would be a
narrow posterior asserting a precision the data does not contain.
 
Because the physical equations do not permit direct evaluation of the likelihood
$p(\mathbf{x} \mid \theta)$, traditional inference methods struggle here. We therefore
use Simulation-Based Inference (SBI) with the BayesFlow library: from a custom
simulator we generate many synthetic pairs $(\theta, \mathbf{x})$ and train a neural
network as an amortized posterior estimator, so that after training the posterior for
any new observation is obtained in a single forward pass, with no re-fitting or MCMC.
 
The primary objective of our analysis is to investigate how the \emph{timing} of the
measurement affects our ability to infer the past. Our hypothesis is that if the model
genuinely captures the physics of chaos, the posterior should widen (lower contraction)
as the observation window occurs later; and, crucially, even as the model becomes less
certain, it should remain well calibrated, with its reported uncertainty faithfully
reflecting the loss of information.

\section{Data}

Because our problem relies on Simulation-Based Inference, we did not collect real-world
experimental data. Our dataset consists entirely of synthetic simulations from a custom
physics engine written in Python. The classical three-body problem has been studied
extensively in dynamical astronomy~\cite{valtonen2006}, and simulation-based inference
with neural networks has recently been applied to a range of inverse problems where the
likelihood is intractable~\cite{radev2023}; our contribution is to use this framework to
quantify how identifiability degrades with observation lateness in a chaotic setting.
 
Since the network learns strictly from generated data, choosing a physically realistic
prior $p(\theta)$ was one of the most critical steps. We restrict to a two-dimensional
plane and fix the three masses to $m_1 = 1/6$, $m_2 = 2/6$, $m_3 = 3/6$ (ratio $1:2:3$,
summing to one). Distinct masses are used deliberately: equal masses would introduce a
permutation symmetry that produces posterior multimodality unrelated to chaos. All
quantities are non-dimensionalised, with $G=1$ and the length unit chosen so the mean
initial separation equals one.
 
The full state normally needs 12 variables (positions and velocities of all three
bodies). Fixing the centre of mass at the origin and setting total momentum to zero lets
us compute the third body from the first two:
\begin{equation}
    \mathbf{r}_3 = -\frac{m_1\mathbf{r}_1 + m_2\mathbf{r}_2}{m_3}, \qquad
    \mathbf{v}_3 = -\frac{m_1\mathbf{v}_1 + m_2\mathbf{v}_2}{m_3}.
\end{equation}
This reduces the target vector $\theta$ to 8 unknowns: the initial positions and
velocities of bodies 1 and 2.
 
\subsection{Sampling the Prior}
We sample each initial condition in two steps. \textbf{Positions:} we place the three
bodies uniformly within a disk (uniform in area), transform to the COM frame, and rescale
so the mean pairwise separation is exactly 1. \textbf{Velocities:} assigning velocities
carelessly usually makes the bodies fly apart. To keep the system bound, we compute the
potential energy $W$ of the sampled positions and draw velocities from
$\mathcal{N}(0, \sigma_v^2)$ with $\sigma_v = \sqrt{|W|/2}$. By the virial theorem this
tunes the expected kinetic energy so that the virial ratio $2T/|W|$ is of order one, i.e.
the system is bound but non-trivial.
 
\subsection{Rejection Criteria}
Chaotic gravity frequently produces configurations that stall the numerical solver, so we
pass every candidate through a strict rejection filter, integrating it over a fine time
grid and rejecting it if any of the following occurs:
\begin{itemize}
    \item \textbf{Ejection:} any two bodies drift further apart than an escape radius of $10.0$.
    \item \textbf{Collision:} any two bodies come closer than a collision radius of $0.03$.
    \item \textbf{Energy drift:} the relative total energy drifts by more than $10^{-8}$,
    indicating numerical error.
\end{itemize}
Because of the chaotic dynamics, a large fraction of candidates was rejected---roughly
$69\%$, dominated by early collisions. We kept generating until we had a clean dataset of
$50{,}000$ training, $2{,}000$ validation and $2{,}000$ test simulations. The three splits
use fixed, non-overlapping random-seed ranges, which guarantees no simulation is shared
between them and rules out data leakage.
 
%=====================================================================
\section{Statistical Model and Simulator}
%=====================================================================
Since we cannot write a closed-form likelihood for chaotic interactions, the likelihood is
defined implicitly by our simulator. This section describes the physics that generates true
trajectories and the measurement model that generates noisy observations.
 
\subsection{Physical Equations and Integration}
The simulator computes each body's acceleration from Newtonian gravity. To prevent forces
from diverging during close encounters (which would stall the integrator), we introduce a
Plummer softening length $\epsilon = 0.02$:
\begin{equation}
    \frac{d^2 \mathbf{r}_i}{dt^2} = \sum_{j \neq i} m_j \,
    \frac{\mathbf{r}_j - \mathbf{r}_i}{\left(\lVert \mathbf{r}_j - \mathbf{r}_i \rVert^2 + \epsilon^2\right)^{3/2}}.
\end{equation}
Because of softening, the conserved energy is not the standard Newtonian energy but the
kinetic energy plus the \emph{softened} potential energy:
\begin{equation}
    E = \sum_{i=1}^{3} \tfrac{1}{2} m_i \lVert \mathbf{v}_i \rVert^2
    - \sum_{i=1}^{3}\sum_{j=i+1}^{3}
    \frac{m_i m_j}{\sqrt{\lVert \mathbf{r}_j - \mathbf{r}_i \rVert^2 + \epsilon^2}}.
\end{equation}
The same $\epsilon$ appears in the accelerations, the potential, and the energy check, so
the model is internally self-consistent. Because energy is exactly conserved in the true
dynamics, its relative drift over an integration is an \emph{independent} indicator of
numerical error---the same quantity used in the rejection filter. We integrate with SciPy's
\texttt{solve\_ivp} using the 8th-order adaptive Runge--Kutta method DOP853, with strict
relative ($10^{-10}$) and absolute ($10^{-12}$) tolerances.
 
\subsection{Measurement Model and Noise}
A key part of the model is realistic observation. An observer rarely catches a chaotic
system exactly at $t=0$; they observe it during a window at some later time. We first
estimate the median Lyapunov exponent of the accepted systems, finding $\lambda \approx 0.72$,
and set the observation window length to $T = 3/\lambda \approx 4.16$ time units. Rather than
starting at $t=0$, each simulation draws a random measurement delay
$t_0 \sim \mathcal{U}(0, 8.32)$, i.e. $t_0 \in [0, 6/\lambda]$. The trajectory is recorded as
$K=10$ equally spaced snapshots over the interval $[t_0, t_0+T]$. The state $\theta$ is
always the state at $t=0$; only the delay before observation changes.
 
This delay is central to the project: it lets a single model span easy (early) to hard
(late) measurements, and it makes the question ``how does recoverability change with
lateness?'' answerable as an in-distribution measurement (Section~\ref{sec:results}). An
earlier design used a fixed window and imitated a later start by truncating the observation
at inference time; because that fed the network input shapes it had never seen in training,
we replaced it with the random-$t_0$ design so every case stays in-distribution.
 
Finally, no measurement is perfect. We add independent Gaussian noise to every recorded
position and velocity. Although position and velocity carry different units and theoretically warrant separate noise scales, we assign a uniform baseline scale of $\sigma_x = \sigma_v = 0.01$ for simplicity, applied i.i.d. per snapshot. The result is a noisy $10 \times 12$ observation matrix $\mathbf{x}$.
 
%=====================================================================
\section{Approximator Architecture}
%=====================================================================
To estimate $p(\theta \mid \mathbf{x}, t_0)$ we build an amortized neural posterior estimator
with BayesFlow, consisting of an adapter, a summary network, and an inference network.
 
\subsection{Adapter}
The adapter standardises all variables (subtracting the mean and dividing by the standard
deviation across the batch) and casts them to \texttt{float32} for stable gradients. It groups
the inputs into three roles: the target parameters $\theta$ (inference variables); the
$10\times12$ noisy trajectory $\mathbf{x}$ (summary variables, passed to the summary network);
and the delay $t_0$ (a direct inference condition). Supplying $t_0$ directly ensures the network
knows how long the system has evolved, which is essential for handling the chaotic loss of
information.
 
\subsection{Summary Network}
Because the observation $\mathbf{x}$ is a time series of 10 sequential snapshots, feeding it into
a plain feed-forward network would discard its temporal structure. We therefore use BayesFlow's
\texttt{TimeSeriesNetwork} as the summary network, configured with an output dimension of 64, so
it compresses the $10\times12$ input into a fixed-length summary vector
$h(\mathbf{x}) \in \mathbb{R}^{64}$. This is a time-aware architecture that processes the
snapshots in order, capturing the temporal dynamics of the trajectory without explicit time-step
embeddings. A time-aware network is the appropriate choice here precisely because permuting the
snapshots would destroy the information; a permutation-invariant set network, suitable for
exchangeable data, would be wrong for an ordered trajectory.
 
\subsection{Inference Network}
The inference network is a continuous normalizing flow based on Flow Matching~\cite{lipman2023}, conditioned on the
summary vector $h(\mathbf{x})$ and the delay $t_0$. It learns to transform samples from a standard
multivariate Gaussian into samples from the target posterior. We chose Flow Matching over a
simpler coupling flow because the posteriors here are expected to be highly non-Gaussian, curved
and multimodal, including ring-shaped supports (Section~\ref{sec:results}), which a coupling
architecture restricted to near-elliptical shapes would struggle to represent. Flow Matching
parameterises a continuous vector field with a multi-layer perceptron, an unconstrained form well
suited to such complex topologies.
 
For both networks we used the default BayesFlow hyperparameters. No explicit regularisation (such
as dropout) was added: the large simulated dataset ($50{,}000$ samples) together with the
measurement noise, which acts as data augmentation, was sufficient to prevent severe overfitting,
as confirmed by the validation-loss behaviour (Section~\ref{sec:training}).

%=====================================================================
\section{Training Procedure}
\label{sec:training}
%=====================================================================
Because our simulator is computationally expensive, we use offline (amortized) training: the
dataset is generated once and stored, then the network is trained on it. Online training, which
regenerates data every epoch, would be prohibitive given the cost of high-accuracy integration
plus rejection sampling.
 
\subsection{Optimization and Learning Rate}
We optimise the summary and inference networks jointly. The Flow Matching objective minimises the
squared error between the model's predicted vector field and the optimal-transport directions from
the Gaussian base to the target. We use the AdamW optimiser with an initial learning rate of
$5\times10^{-4}$ and a cosine-decay schedule with warmup: the rate ramps up over the first few
epochs to stabilise early gradients, then anneals smoothly toward zero.
 
\subsection{Training Budget and Early Stopping}
We use a batch size of 64; with $50{,}000$ training samples this is $781$ steps per epoch. We set a
maximum of 150 epochs but monitor the validation loss (on the $2{,}000$-sample validation set) and
apply early stopping with a patience of 20 epochs. The validation loss fell sharply at first, then
plateaued and drifted slightly upward; with no improvement for 20 consecutive epochs, training
halted at epoch 88.
 
Crucially, we did \emph{not} keep the final-epoch weights. We restored the weights from the epoch
with the lowest validation loss, epoch 68. Saving the best-validation model rather than whichever
epoch training happened to stop on was worth roughly a factor of four in calibration error, which is a
larger effect than any architecture or dataset change we tried.
 

 
%=====================================================================
\section{Diagnostics}
%=====================================================================
Before trusting the estimator we validate it on the untouched $2{,}000$-sample test set. Throughout,
we keep two questions separate: \emph{how much} the measurement tells us (contraction) and whether
the stated uncertainty is \emph{honest} (calibration). These are orthogonal, which means a narrow posterior can
be miscalibrated, a wide one is well calibrated.
 
\subsection{Convergence}
The validation loss decreased rapidly over the first ${\sim}40$ epochs before stabilising. Early
stopping and restoration of the epoch-68 weights (validation loss $\approx 1.066$) gave a model that
generalises well and avoids the slight overfitting seen in the final epochs.
 
\subsection{Simulation-Based Calibration (SBC)}
SBC verifies whether the stated uncertainty is honest. For each test simulation we draw 100 posterior
samples and compute the rank of the true $\theta$ among them; if the model is perfectly calibrated,
these ranks are uniformly distributed. We plot the rank ECDF against $95\%$ confidence bands (we use
the ECDF form rather than rank histograms because it admits simultaneous confidence bands and needs no
binning choice). Across the 8 parameters the ECDF curves lie only slightly outside the bands
(see Appendix, Figure~\ref{fig:ecdf}). This small residual is expected when pooling test cases of very different
posterior widths (mixing tight early measurements with wide late ones); the mean calibration error
remains low ($\approx 0.02$), confirming that when the model is uncertain it is honest about it.
 
\subsection{Posterior Contraction}
To measure how much information the observation adds relative to the prior, we compute the posterior
contraction per parameter:
\begin{equation}
    c = 1 - \frac{\mathrm{Var}[\theta \mid \mathbf{x}]}{\mathrm{Var}[\theta]}.
\end{equation}
A value near 1 means a tight, informative posterior; near 0 means the observation added almost
nothing. Pooled over the whole test set the contraction is moderate and varies by parameter
(see Appendix, Figure~\ref{fig:contraction}). But a pooled value is of limited use in a chaotic system. Its real
value lies in how it \emph{fades} with the measurement delay $t_0$, which we examine next.
 
\subsection{Recovery}
As an additional check we plot the posterior mean against the ground truth for each parameter
(see Appendix, Figure~\ref{fig:recovery}), with correlations in the range $0.5$--$0.7$. The scatter is substantial,
but this is expected rather than a defect: perfect accuracy is not the goal, and for hard (late) cases
the posterior mean should \emph{not} coincide with the true value, since the observation does not
determine it. The rigorous test of correctness is calibration, not point accuracy.

%=====================================================================
\section{Inference and Results}
\label{sec:results}
%=====================================================================
Because the project is entirely simulation-based, we have no external real-world dataset; we perform
inference on the held-out test set, which provides both the observation and the ground-truth initial
condition.
 
The core question is whether the network correctly learns that chaos destroys information over time.
To answer it, we bin the $2{,}000$ test simulations by their measurement delay $t_0$ (in Lyapunov
units, $t_0 \times \lambda$) and compute the mean contraction and calibration error per bin. Because
the model was trained on the same $t_0$ range, every bin is in-distribution, so this reads how
identifiability changes with lateness honestly.
 
\subsection{Chaos Summary}
The resulting summary (see Appendix, Figure~\ref{fig:chaos}) shows a clear trend consistent with physical theory.
For early measurements ($t_0 \times \lambda \approx 0.5$), the mean contraction is strong, around
$0.76$: chaos has not yet scrambled the system, so the network narrows down the initial conditions
well. As the delay grows, contraction falls steadily, reaching about $0.26$ for late measurements
($t_0 \times \lambda \approx 5.5$)---the posterior has widened until it is nearly as broad as the
prior. Yet across all delays the calibration error stays flat, between $0.02$ and $0.03$. In short:
\emph{contraction fades, calibration survives}. As information about the initial condition decays
exponentially, the model becomes appropriately less certain while remaining honest about exactly how
much less certain it is.
 
\subsection{Posterior Predictive Checks and Corner Plots}
To check that these wide posteriors are physically realistic rather than a mathematical artifact, we perform a posterior predictive check on a late-measurement case: we draw samples from the posterior and re-simulate them forward. Within the observation window $[t_0, t_0+T]$, every re-simulated path closely tracks the observed noisy points---the samples are consistent with the data. But traced back
into the unobserved past $[0, t_0]$, the paths spread vastly apart, covering roughly 7 units of space
(see Appendix, Figure~\ref{fig:ppc}). This is the visual signature of chaos: completely different histories are
compatible with the same late measurement. We stress that this check establishes consistency with the
data, not closeness to the truth; only the calibration analysis establishes correctness.
 
Finally, the full corner plot for a late measurement (see Appendix, Figure~\ref{fig:corner}) reveals the value of the
Flow Matching architecture. Rather than an elliptical Gaussian blob, the posterior for the position of
body 1 forms a distinct \emph{ring}: the distance of the body from the origin is well constrained, but
its angular direction has been erased by chaotic mixing. A method restricted to elliptical posteriors
could not represent this; the continuous normalizing flow captures it naturally.
 
%=====================================================================
\section{Limitations and Potential Improvements}
%=====================================================================
Two limitations are worth stating explicitly.
 
\paragraph{Pooled calibration.} Our SBC rank ECDFs lie slightly outside the $95\%$ band when evaluated
over the entire test set. Within each $t_0$ bin the calibration error is low ($\approx 0.02$); the
pooled test, however, has far more statistical power and detects a small residual miscalibration that
the per-bin measure cannot. Because the pooled set mixes tight early posteriors with wide late ones, the
global metric can appear slightly distorted even when the model is well calibrated at any specific delay.
A natural next step is to run a full SBC on narrow $t_0$ slices to confirm the residual is uniform across
delays.
 
\paragraph{Posterior support beyond the prior.} The Flow Matching network has no built-in knowledge of
physical constraints. Prior samples are filtered by rejection, but posterior samples are not, and the
smooth flow never saw the rejected region. As a result a small fraction of posterior draws (about
$1.7\%$) land on near-collision configurations that the prior would have discarded, more often for late
measurements where the posterior is wider. The interpretation is that the one-dimensional marginals are
well calibrated but the joint distribution is not perfectly so: the set of initial conditions compatible
with a measurement is thin and filament-like, and a smooth flow reproduces the marginal widths but not
that fine geometry. Two improvements suggest themselves: embedding conservation laws into the loss (a
physics-informed approach) to discourage impossible draws; and running SBC on conserved quantities such
as energy and angular momentum, nonlinear functions of all eight parameters that would probe the joint
structure directly, and which we already record but have not yet used as calibration test quantities.
 
%=====================================================================
\section{Conclusion}
%=====================================================================
In this project we applied Simulation-Based Inference to the classical chaotic three-body problem.
Traditional methods often fail on chaotic systems because they seek a single, overconfident estimate,
ignoring that chaos destroys historical information.
 
Using an amortized neural posterior estimator based on Flow Matching, and explicitly feeding the
measurement delay $t_0$ to the network, we quantified this loss of information. Three points summarise
the work. First, in a chaotic system a wide posterior is the correct answer, not a failure: information
decays exponentially at a rate set by $\lambda \approx 0.72$, and a method that stayed confident as the
measurement started later would merely be misreporting its uncertainty. Second, contraction fades but
calibration survives: across delays from $0.5$ to $5.5$ Lyapunov times, contraction falls from $0.76$ to
$0.26$ while the calibration error stays flat between $0.02$ and $0.03$. Third, amortization is what made
the analysis possible: a single model, trained once with $t_0$ as a known input, spans every delay
regime, letting us measure how difficulty grows with $t_0$ over the whole test set at negligible cost.
In a chaotic universe, being uncertain is scientifically accurate, and reporting that uncertainty
correctly is the mark of an honest statistical model.
 
%=====================================================================
\section{Reflection on Own Learnings}
%=====================================================================
The most valuable lesson of this project was learning to treat a wide posterior as a result rather than
a failure. Our intuition at the start was that a good model should produce sharp, confident answers, and
an early design did exactly that, but its calibration checks failed systematically, exposing it as
overconfident. Seeing this concretely taught us the difference between \emph{knowing a lot} and
\emph{honestly reporting how much one knows}, and why calibration, not sharpness, is the right criterion
for a posterior in a chaotic system.
 
We also learned how much of simulation-based inference lives in the modelling choices around the network
rather than in the network itself: designing a proper prior, keeping the simulator self-consistent (for
instance using the same softening length in both the force and the energy), shaping the observation model
so that a scientific question becomes an in-distribution measurement, and validating with SBC and
posterior contraction. On the engineering side, the single most impactful change was saving the
best-validation model rather than the final one, which is exactly a reminder that careful methodology can matter more
than model complexity.


\newpage
\begin{thebibliography}{9}

\bibitem{valtonen2006}
Valtonen, M., \& Karttunen, H. (2006). \emph{The Three-Body Problem.} Cambridge University Press.

\bibitem{radev2023}
Radev, S.~T., Schmitt, M., Schumacher, L., Elsem\"uller, L., Pratz, V., Sch\"alte, Y., K\"othe, U.,
\& B\"urkner, P.-C. (2023). \emph{BayesFlow: Amortized Bayesian Workflows with Neural Networks.}
Journal of Open Source Software, 8(89), 5702.
 

 
\bibitem{lipman2023}
Lipman, Y., Chen, R.~T.~Q., Ben-Hamu, H., Nickel, M., \& Le, M. (2023).
\emph{Flow Matching for Generative Modeling.} International Conference on Learning Representations (ICLR).
\end{thebibliography}

\newpage
\appendix
\section{Appendix: Figures}

\begin{figure}[htbp]
    \centering
    \includegraphics[width=0.8\textwidth]{example_orbits.png}
    \caption{Three simulated three-body systems. Identical physics, three different initial conditions, three completely different trajectories.}
    \label{fig:illustrative}
\end{figure}

\begin{figure}[htbp]
    \centering
    \includegraphics[width=\linewidth]{calibration_ecdf.png}
    \caption{Rank ECDF deviation from uniform for each parameter, pooled over all $t_0$, with
    $95\%$ confidence bands.}
    \label{fig:ecdf}
\end{figure}

\begin{figure}[htbp]
    \centering
    \includegraphics[width=0.65\linewidth]{contraction.png}
    \caption{Posterior contraction per parameter, pooled over all $t_0$. Moderate values are
    expected because the test set mixes easy and hard cases.}
    \label{fig:contraction}
\end{figure}

\begin{figure}[htbp]
    \centering
    \includegraphics[width=\linewidth]{recovery.png}
    \caption{Posterior mean against ground truth for all test simulations, with credible intervals.
    Scatter is expected; accuracy is not the goal.}
    \label{fig:recovery}
\end{figure}

\begin{figure}[htbp]
    \centering
    \includegraphics[width=0.8\linewidth]{chaos_summary.png}
    \caption{Contraction (blue, left axis) and calibration error (red, right axis) versus measurement
    lateness. Contraction falls from $0.76$ to $0.26$ while calibration error stays flat at
    $0.02$--$0.03$. This is the central result of the project.}
    \label{fig:chaos}
\end{figure}

\begin{figure}[htbp]
    \centering
    \includegraphics[width=\linewidth]{posterior_predictive.png}
    \caption{Re-simulated trajectories versus the noisy observation, for an early (left) and a late
    (right) measurement. Solid: the observation window. Faint dashed: the earlier, unmeasured interval.}
    \label{fig:ppc}
\end{figure}

\begin{figure}[htbp]
    \centering
    \includegraphics[width=0.6\linewidth]{posterior_example.png}
    \caption{Posterior samples for the initial position of body 1. Blue: measured early. Red: measured
    late. The cross marks the true value. The late case forms a ring.}
    \label{fig:corner}
\end{figure}

\end{document}